\section{Advance by a bounded chunk, including its padding}
Each ordinary chunk has a four-byte identifier, a four-byte little-endian payload
length and that many payload bytes. An odd payload length has one following pad byte.
The length does not include that pad or the eight-byte chunk header. Even an empty chunk advances past its header. Compare all four
identifier bytes: the format identifier is \texttt{fmt\ } with a final space.

At a candidate header offset, require at least eight remaining bytes before loading
its fields. After consuming that header, compare the payload length plus its parity
bit with the remaining extent. Use unsigned wide arithmetic: a declared length of
\texttt{ffffffff} needs 4294967296 bytes including its pad, not zero. Reject it against
the finite remaining interval before advancing an address.

The first ancillary chunk in the worked file has identifier \texttt{JUNK}, length
three, payload \texttt{01 02 03} and pad \texttt{a5}. It occupies twelve bytes in total,
so the following format header begins at 24. Treat the pad as an extent byte, not as
a fourth payload byte. This profile permits any pad value. After skipping an unknown
chunk, resume at the computed next header; do not search its payload for \texttt{fmt}
or \texttt{data} text that might occur by chance.
